\scanpage{55}
\mspage{\handwritten{- 10 -}}
% archive sheet number: 54
\uline{sa somme est une fonction holomorphe dans ce polycylindre}.

De même, en vertu du cor. du th.2 :

PROPOSITION 6.- \uline{Soit} $B$ \uline{un ensemble ouvert dans} ${K}^{q}$, $g_k$ $(1\handwritten{\leq}k\handwritten{\leq}p)$ $p$ \uline{fonctions scalaires holomorphes dans} $B$, $A$ \uline{un ensemble ouvert dans} $K^{p}$ \uline{contenant l'image de} $B$ \uline{par} $(g_1,\ldots,g_p)$. \uline{Si} $\handwritten{f}$ \uline{est une fonction holomorphe dans} $A$, \uline{la fonction composée} $\handwritten{f}(g_1,\ldots,g_p)$ \uline{est holomorphe dans} $B$.

De façon imagée, toute fonction \uline{composée} de fonctions holomorphes est holomorphe.

PROPOSITION 6 bis.- \uline{Soient} $E_1,\ldots,E_r$ \uline{et} $F$ $r+1$ \uline{espaces de Banach sur} $K$, \uline{et soit} $(z_1,\ldotp\ldotp,z_r)\handwritten{\longrightarrow}\handwritten{[}z_1.z_2\ldots z_r\handwritten{]}$ \uline{une application multilinéaire continue de} $E_1\handwritten{\times}\ldots\handwritten{\times}E_r$ \uline{dans} $F$. \uline{Si} $\handwritten{f}_1,\ldots,\handwritten{f}_r$ \uline{sont} $r$ \uline{fonctions holomorphes dans un ensemble ouvert} $A\handwritten{\subset}K^p$, \uline{à valeurs dans} $E_1,\ldots,E_r$ \uline{respectivement}, \uline{la fonction} $\handwritten{[}\handwritten{f}_1.\handwritten{f}_2\ldots\handwritten{f}_r\handwritten{]}$ \uline{est une fonction holomorphe dans} $A$, \uline{à valeurs dans} $F$.

% The right edge clips the first proof line after the visible letters remarqu.
Bornons-nous pour simplifier au cas $r=2$. Il suffit alors de remarqu que si $(x_{\handwritten{\lambda}})_{\handwritten{\lambda\in}L}$ est une famille absolument sommable dans $E_1$ de somme $S$, et $(\handwritten{y}_{\handwritten{\mu}})_{\handwritten{\mu\in}M}$ une famille absolument sommable dans $E_2$, de somme $\handwritten{t}$, la famille $(\handwritten{[}x_{\handwritten{\lambda}}.\handwritten{y}_{\handwritten{\mu}}\handwritten{]})_{(\handwritten{\lambda},\handwritten{\mu})\handwritten{\in}L\handwritten{\times}M}$ est absolument sommable dans $F$ et a pour somme $\handwritten{[}\handwritten{s}.\handwritten{t}\handwritten{]}$, puisqu'il existe\handwritten{\unsure{.}}une constante $k\handwritten{>}0$ telle que $\handwritten{\|[}z_{\handwritten{1}}.z_{\handwritten{2}}\handwritten{]\|}\handwritten{\leq}k.\handwritten{\|}z_{\handwritten{1}}\handwritten{\|}.\handwritten{\|}z_{\handwritten{2}}\handwritten{\|}$ ; on appliquera cette remarque aux séries entières dont les sommes sont égales à $\handwritten{f}_1$ et $\handwritten{f}_2$ au voisinage d'un point de $A$.

COROLLAIRE.- \uline{Soit} $\handwritten{f}=(f_1,\ldots,f_q)$ \uline{une fonction définie dans un ensemble ouvert} $A\handwritten{\subset}K^p$ \uline{et à valeurs dans} $K^q$. \uline{Pour que} $\handwritten{f}$ \uline{soit holomorphe dans} $A$, \uline{il faut et il suffit que chacune des fonctions scalaires} $f_k$ $(1\handwritten{\leq}k\handwritten{\leq}q)$ \uline{le soit}.
